A performance measure is declared in the adjoint input file with one of three forms. The \emph{direct} form accumulates a quantity, the \emph{residual} form accumulates the squared difference between a quantity and an observed counterpart, and the \emph{instantaneous} form averages a quantity over the times at which it is observed. This chapter describes the third.

\section*{Motivation}

The direct form answers a question about a total: how much water left through a set of boundaries over the simulation, and how does that total respond to a parameter. Written over the time steps $k$ at which the measure has entries,

\begin{equation}
	\label{eqn:direct}
	J_{\mathrm{direct}} = \sum_{k} F^{k}.
\end{equation}

That total depends on how the simulation is discretized in time. Splitting a stress period into more time steps adds terms to the sum. When the question is instead about a state at a particular time, such as the head at an observation well or the rate a stream is losing at the end of a season, a sum over time steps is not the quantity of interest and its value moves with the time discretization.

The instantaneous form is the time-weighted mean

\begin{equation}
	\label{eqn:instantaneous}
	J_{\mathrm{inst}} = \frac{\sum_{k} \Delta t^{k} F^{k}}
	{\sum_{k} \Delta t^{k}},
\end{equation}

\noindent where $\Delta t^{k}$ is the length of time step $k$ and the sums run only over the time steps at which the measure has entries. Refining the time discretization leaves equation~\ref{eqn:instantaneous} unchanged, because both sums scale together.

\section*{Implementation}

Two things follow from equation~\ref{eqn:instantaneous}, and both differ from the direct and residual forms.

\subsection*{Time steps are independent}

A direct or residual measure accumulates through time, so the adjoint solution at one time step feeds the step before it. That coupling is carried by the storage term $\partial \mathbf{b}^{k} / \partial \mathbf{h}^{k-1}$ described in Chapter~\ref{ch:storage}: the right-hand side of the adjoint problem at step $k$ is

\begin{equation}
	\label{eqn:carryback}
	\frac{\partial \mathbf{b}^{k+1}}{\partial \mathbf{h}^{k}}
	\bm{\lambda}^{k+1}
	- \frac{\partial F^{k}}{\partial \mathbf{h}^{k}},
\end{equation}

\noindent which propagates the later step's adjoint state backward in time.

An instantaneous measure asks a separate question at each observation time, so no such propagation applies and the first term of equation~\ref{eqn:carryback} is dropped. Each observed time step is solved on its own and the results are combined by equation~\ref{eqn:instantaneous}. The same reasoning removes the carry-back of a lake stage, described in Chapter~\ref{ch:lake}.

\subsection*{Unobserved time steps are skipped}

Because the steps are independent, a time step at which the measure has no entries contributes nothing and is not solved at all. For a measure observed at one time in a long simulation this is the difference between one adjoint solve and one per time step.

\section*{On examples}

The chapters that follow close with a numerical comparison, because in each case a sensitivity can be checked against a finite-difference derivative and the agreement is the evidence that the formulation is right. No comparable example is given here. The instantaneous form does not change any single sensitivity; it changes which time steps are solved and how their results are combined. A worked example would show the same per-step sensitivities as the direct form, rescaled by a constant, which demonstrates nothing that equation~\ref{eqn:instantaneous} does not already state.

The property worth testing is the one that motivates the form: that the result does not depend on the time discretization. Halving the time step length while holding the simulated period fixed leaves an instantaneous measure unchanged and doubles the number of terms in a direct measure.
